\chapter*{Abstract}

\section*{\tituloPortadaEngVal}

The study of large musical corpora constitutes a fundamental tool in disciplines such as historical musicology, ethnomusicology, and musical analysis. Traditionally, this type of research has depended on manual processes of score reading and cataloging, a rigorous methodology but limited in terms of scalability against the growing volume of digitized musical material. Despite the emergence of standardized symbolic formats like MusicXML or MEI, critical challenges persist regarding the heterogeneity of the corpora, inconsistencies in symbolic standards, algorithmic complexity in detecting rhythmic phenomena, and the division between structural analysis and the semantic analysis of lyrics. Furthermore, general-purpose language models (LLMs) present severe limitations in context and reasoning when executing specialized musical analysis in real time over extensive collections.

In this context, this work presents the design and implementation of a hybrid architecture that transforms a heterogeneous corpus of musical scores into a structured set of musicological variables through automated extraction pipelines (based on music21), storing them in a relational SQLite database that is then queried using a local language model through SQL query generation (Text-to-SQL) and an MCP server. The experimental results indicate that, in several specialized musicological tasks, the proposed architecture achieves higher accuracy than state-of-the-art commercial models such as Gemini 2.5 Flash when these systems rely exclusively on basic document-retrieval-based Retrieval-Augmented Generation (RAG) mechanisms.

\section*{Keywords}

\noindent Computational musicology, Musical corpora, Feature extraction, Language models (LLMs), Music21, Folk music, Information retrieval, Model Context Protocol (MCP), Text-to-SQL, Digital Humanities
